\documentclass[10pt,a4paper]{amsart}
\usepackage[margin=19mm]{geometry}
\usepackage{amsmath,amssymb,mathtools}
\usepackage[scr=rsfs,cal=euler]{mathalfa}
\usepackage{xfrac}
\usepackage[unicode,hidelinks,bookmarks=false]{hyperref}
\newcommand{\ie}{i.e.\ }
\newcommand{\cf}{cf.\ }
\newcommand{\resp}{respectively\ }
\newcommand{\Subsection}{\subsection}
\providecommand{\nobreakdash}{\nobreak\hbox{-}}
\setlength{\emergencystretch}{2em}
\multlinegap0pt
\usepackage{bm}
\usepackage{bbm}
\usepackage{dsfont}
\usepackage{xspace}
\usepackage{cancel}
\usepackage{mathtools}
\usepackage{amsthm}
\makeatletter
\@ifundefined{lemma}{\newtheorem{lemma}{Lemma}}{}
\makeatother
\def\ERns{Erd\"os-R\'enyi}
\title[Unusual properties of contact processes on percolated graphs]{Unusual properties of contact processes on percolated graphs}
\author{Rick Durrett}
\date{Source: arXiv:2403.18592v2 (CC BY 4.0)}
\begin{document}
\maketitle

\noindent\emph{Source and licence.} Unusual properties of contact processes on percolated graphs. Version arXiv:2403.18592v2 (CC BY 4.0), distributed under the Creative Commons Attribution 4.0 International licence, as recorded in the arXiv licence metadata for this exact version and shown on the abstract page https://arxiv.org/abs/2403.18592v2. Full rights notice: licenses/arxiv-2403.18592-CC-BY-4.0.txt.\par\smallskip

\noindent\textit{Editorial scope.} Counted unit: the Lemma whose source label is \texttt{maxpER} in the article (source0.tex lines 523--526, source label \texttt{maxpER}), delivered together with the article's own complete original proof of it (source0.tex lines 528--557, one \texttt{proof} environment of 30 lines). No mathematics is written, repaired or re-derived: statement and proof are the article's own text line for line. Required context retained uncounted: none. Every definition, display and label of those lines is retained unchanged. Omitted from this package and not redistributed: 1--522, 527--527, 558--893. Nothing inside a retained excerpt is omitted or altered: every retained body is delivered exactly as the article's lines are written, so no omission receipt of any kind accompanies this package. Citations: the retained excerpts carry no citation command at all, so no bibliography entry of the article is transcribed and no reference is redistributed. Reference adaptation: none was needed and none was made; every reference inside the delivered unit resolves inside it, so no unresolved reference or citation is produced. Printed slips kept literal: no printed word, symbol, hypothesis or number of the counted statement or of its proof was altered. Layout and class: the article's own class is \texttt{article} with packages including latexsym, amssymb, amsmath, amsfonts, amsthm, graphicx, url; the wrapper is the lane's pinned \texttt{amsart} editorial wrapper at 10pt on A4, which restates the theorem-like environments the retained text needs and adds the article's own macro definitions that the retained text needs (\texttt{\textbackslash ERns}), with extra wrapper packages (bm, bbm, dsfont, xspace, cancel, mathtools). The delivered numbering is the unit's own. Assets: no retained span contains a figure environment, a graphics-inclusion command, a PDF-page inclusion, a table or a TikZ picture; no image file is copied into this package, the package declares no asset dependency and no image file is redistributed with it. Licence: version arXiv:2403.18592v2 of this article is distributed under the Creative Commons Attribution 4.0 International licence, as recorded in the arXiv licence metadata for this exact version and shown on the abstract page https://arxiv.org/abs/2403.18592v2. A full rights notice is delivered at licenses/arxiv-2403.18592-CC-BY-4.0.txt. Characters: every retained excerpt is pure ASCII, so the delivered engine, XeLaTeX with its default Latin Modern text font, reports no missing character.
\begin{lemma} \label{maxpER}
If $\nu<1$ then as $N \to\infty$ the longest path in an \ERns$(N,\nu/N)$
$\sim (\log N)/\log(1/\nu)$.
\end{lemma}

\begin{proof}
If $k^2/N \to 0$, which we assume throughout the proof, then the expected number of (self-avoiding) paths of length $k$ in an Erd\"os-R\'eny($N,\nu/N$) graph
$$
\Pi(N,k,\nu) = N(N-1) \cdots (N-k+1) \cdot (\nu/N)^{k-1} \sim N \nu^{k-1} .
$$
$\Pi(N,k_1,\nu) \approx 1$ when $k_1(\nu,N)= (\log N)/\log (1/\nu)$. If $\bar k = (1+\delta)k_1$ then 
$$
\Pi(N,\bar k,\nu) \le N^{-\delta} \to 0 ,
$$
which gives the upper bound on the length of the longest path. 

If $\underline{k} = (1-\delta) k_1$ then $\Pi(N,\underline{k},\nu) \sim N^\delta$. To prove the lower bound let $\pi$ be a sequence of distinct vertices of length $\underline{k}$, let $A_\pi$ be the event that $\pi$ is a paths of length $\underline{k}$ in the graph, and let $Y_\delta$ be the number of path of length $\underline{k}$. If $\pi$ and $\sigma$ do not share an edge in common then $A_{\pi}$ and $A_\sigma$ are independent. Let $\Sigma^j$ be the number of pairs of paths $\pi$ and $\sigma$ that have exactly $j$ edges in common. 

We get a lower bound on $\Sigma^0$ by assuming all vertices in each path are different.
$$
N(N-1) \cdots (N-2k-1)  (\nu/n)^{2k-2} \le \Sigma^0  \le \Pi^2 ,
$$
so $\Sigma^0 \sim \Pi^2= (N \nu^{k-1})^2$. In $\Sigma^1$ we need to pick the edge to be the same.
\begin{align*}
\Sigma^1  \le N^k \cdot k N^{k-2} \cdot ( \nu/N)^{(k-1)+(k-2)}& \\
\sim kN\nu^{2k-3}= (\Pi)^2 \cdot k \nu^{-1}/n = o(\Sigma^0). &
\end{align*}
 When it comes to $\Sigma^2$ the two edges can be adjacent in the path or not
\begin{align*}
\Sigma^2 \approx N^k \cdot [k N^{k-3} + k^2 N^{k-4}] 
\cdot ( \nu/N^{(k-1)+(k-3)})& \\
\le (\Pi)^2 \cdot \nu^{-2}  [k/N+ k^2/N^2] = o(\Sigma^0)&
\end{align*}
The number of possibilities increases as the number of duplicates increases but the best case occurs when all the agreements are in a row. Since $\sigma^0 \sim \Pi^2$ the square of the mean, it follows that the variance is $o(\Pi^2)$, and the ratio $Y_\delta/EY_\delta \to 1$.
\end{proof}

\end{document}
